\section{Implementation Details}

\subsection{System Components}

The system is implemented using a modular architecture, where each component is responsible for a specific function in the query processing pipeline.

The main components include:
\begin{itemize}
    \item \textbf{Router module:} Determines the execution path for each query.
    \item \textbf{Workflow manager:} Orchestrates the end-to-end processing pipeline.
    \item \textbf{Retrieval module:} Handles tenant-scoped vector-based similarity search.
    \item \textbf{LLM service:} Manages inference using the shared base model.
    \item \textbf{Memory store:} Maintains tenant-specific data and context.
    \item \textbf{Tool module:} Processes structured data queries (e.g., Excel, CSV).
\end{itemize}

This modular design improves flexibility, scalability, and maintainability of the system.

\subsection{Multi-tenant Design}

Each tenant in the system is associated with:
\begin{itemize}
    \item A unique tenant identifier
    \item An isolated knowledge base
    \item An isolated memory context
    \item An optional LoRA personalization adapter
\end{itemize}

The system ensures strict isolation between tenants by:
\begin{itemize}
    \item Separating vector database partitions
    \item Maintaining independent memory contexts
    \item Routing queries based on tenant identifiers
\end{itemize}

At the same time, all tenants share a single base model, enabling efficient resource utilization.

\subsection{Routing Optimization}

The routing mechanism is designed to reduce unnecessary computational overhead and to better align execution cost with query requirements.

Instead of processing all queries through the LLM, the system dynamically selects the most efficient execution path:

\begin{itemize}
    \item \textbf{Tool-based processing:} For structured data queries such as Excel or CSV
    \item \textbf{Retrieval-based generation:} For document-based queries (e.g., PDF)
    \item \textbf{Direct LLM inference:} For general queries
\end{itemize}

This approach reduces inference latency and token usage by avoiding a uniform pipeline for all requests.

\subsection{Resource Optimization Strategy}

The system incorporates multiple techniques to optimize resource usage:

\begin{itemize}
    \item \textbf{Shared model architecture:} A single base LLM serves all tenants
    \item \textbf{Dynamic adapter loading:} Tenant-specific LoRA adapters are loaded on demand
    \item \textbf{Selective LLM invocation:} Retrieval and full generation are used only when justified by the route decision
\end{itemize}

These strategies reduce memory footprint and improve scalability in cloud environments.
